\documentclass[11pt]{article}

\usepackage[margin=1in]{geometry}
\usepackage{amsmath,amssymb,array,booktabs,lmodern,xspace}
\usepackage[T1]{fontenc}
\usepackage[round]{natbib}
\usepackage[hidelinks]{hyperref}

\newcommand{\pkg}[1]{{\normalfont\fontseries{b}\selectfont #1}}
\newcommand{\proglang}[1]{\textsf{#1}}
\newcommand{\code}[1]{\texttt{\detokenize{#1}}}
\newcommand{\citrees}{\pkg{citrees}\xspace}
\newcommand{\sklearn}{\pkg{scikit-learn}\xspace}

\renewcommand{\thesection}{S\arabic{section}}
\renewcommand{\thetable}{S\arabic{table}}
\renewcommand{\theequation}{S\arabic{equation}}

\title{Supplement to ``\citrees: Conditional Inference Trees and Forests for
  \proglang{Python}''}
\author{Robert Milletich \and Justin Downes \and Newel Hirst}
\date{}

\begin{document}

\maketitle

This supplement holds the full designs and the secondary results of the
validation studies in the article. Section and table numbers without the
prefix S refer to the article. Every number here is regenerated by the
replication command of Section~9, and Section~\ref{sec:producer-map} maps each
reported number to the component and output file that produce it.

\section{Selector projection sensitivity} \label{sec:rdc-sensitivity}

The RDC implementation uses 10 random projections, which give 20 sine and
cosine columns per variable. A sensitivity analysis with 5 seeds and 5 folds
compared 5, 10, 20, and 40 projections on Wine (178 samples, 13 features),
Glass (214 samples, 10 features), Heart Statlog (270 samples, 13 features),
and Parkinsons (195 samples, 22 features). Each ranking was evaluated by
fitting three downstream classifiers (logistic regression, a support vector
machine, and $k$-nearest neighbors) to its top 5 and top 10 features, or top
5 only for Glass, which has 10 features.

Across 25 paired seed and fold evaluations per dataset, the median selector
time with 20 projections was 3.00 to 3.06 times that with 10. At these
feature counts, the mean balanced accuracy difference, 20 minus 10
projections, ranged from $-0.0052$ to $0.0005$ across the three downstream
classifiers. The mean overlap between the selected feature sets ranged from
95.4\% to 99.8\%. Ten projections therefore cut the median selector time by
about two thirds relative to 20, with no measurable loss.

\section{Design of the cardinality and matched studies} \label{sec:cardinality-design}

\subsection{Cardinality study}

Each replicate has 200 observations and five independent ordered discrete
predictors with exactly 2, 4, 10, 20, and 200 observed levels. Every level
occurs equally often within a predictor. Values within columns and the five
column positions are shuffled independently in each replicate, and the
response is generated independently of all predictors.

Tree comparisons use 10,000 replicates for classification and for
regression. \citrees and \pkg{partykit} use a familywise level of 0.05 with
Bonferroni adjustment across the five predictors and a nominal budget of 999
permutations. \citrees evaluates each candidate with 4,995 permutations to
keep its resolution at the adjusted level, and \pkg{partykit} uses its native
Monte Carlo distribution with 999 permutations. CART stumps are the
impurity-based comparison. The analysis records candidate statistics,
$p$-values, native root decisions, ties, and the no-split outcome. The
\pkg{partykit} native root decisions split in 4.05\% of classification and
4.98\% of regression replicates.

Forest comparisons use 2,500 replicates and 100 depth-one trees per forest.
Every tree receives all five predictors and a bootstrap sample of 200
observations. The comparison includes \citrees forests, \pkg{partykit}
conditional inference forests, and \sklearn random forests. Each
implementation contributes its documented native summary: normalized impurity
decrease for \citrees and \sklearn, and unconditional permutation importance
with 10 permutations for \pkg{partykit}. Positive associations between
cardinality and selection or importance are tested with 9,999 joint label
randomizations and the add-one rule. Holm adjustment is applied to the
prespecified family of 14 tree and forest trend tests of Table~6. These tests
are upper-tailed, so they do not detect a preference for low-cardinality
predictors.

\subsection{Matched comparison with partykit}

The matched comparison partitions each dataset by 10 repeats of five-fold
cross-validation. Trees and forests from \citrees and \pkg{partykit} receive
identical training and evaluation rows and the same integer seed through
their native random-number streams. The controls align the data partitions,
stopping depth, node-size constraints, predictor availability, permutation
budget, bootstrap size, and forest size. The primary comparison runs \citrees
with feature and threshold scanning off, so candidates are tested in random
order. A scanning control repeats the tree comparison with scanning on.

Tree comparisons record whether each implementation splits and whether the
root predictors agree when both split. Tree summaries count accepted splits.
Forest summaries use normalized impurity decrease for \citrees and
unconditional permutation importance with 10 permutations for \pkg{partykit}.
When both summaries distinguish among predictors, concordance is measured by
Kendall's $\tau_b$ and by the Jaccard similarity of the top five sets,
including ties at the fifth position. Held-out classification comparisons
report balanced accuracy, prediction agreement, probability differences and
rank correlation, area under the ROC curve, and log loss. Regression
comparisons report root mean squared error, mean absolute error, prediction
correlation, and the mean absolute difference between predictions. Every fit
is repeated with the same seed. Estimates first average folds within each
partition repeat and then average the 10 repeat means, and percentile
intervals use 9,999 bootstrap samples of the repeat means. The design
contains 200 paired dataset, partition, and model-family comparisons.

The following results complete Section~7.2 of the article. With scanning
off, root-predictor agreement was 8\% in classification (95\% interval 2\% to
14\%) and 12\% in regression (4\% to 20\%). Tree-level summary concordance
was 0.153 in classification and 0.698 in regression. The bootstrap interval
for the difference in forest balanced accuracy, 0.002 to 0.008, excludes zero
but is small on the accuracy scale. Tree-level predictions agreed on 91.4\%
of classification test cases and correlated at 0.789 in regression.

\section{Synthetic scaling of forest fitting time} \label{sec:scaling}

Table~\ref{tab:performance-scaling} extends the real-data comparison of
Table~8 to synthetic Gaussian predictors, the continuous case, as the sample
size grows. The benchmark configuration is 1.2 to 3.2 times faster than
\pkg{partykit} on one core at every sample size and 6.6 to 16 times slower
than \pkg{partykit} on 32 cores. On these predictors the max-type test is 13
to 36 times faster than the per-threshold Bonferroni test and below the
\pkg{partykit} 32-core time at every sample size. A strong-signal variant
changes the benchmark times by at most 19\% except at 8,000 observations,
where it takes 1.6 to 1.7 times longer, and changes the \pkg{partykit} times
by at most 18\%.

\begin{table}[t!]
\centering
\small
\setlength{\tabcolsep}{4pt}
\begin{tabular}{rrrrrrr}
\toprule
 & & \multicolumn{3}{c}{\citrees} & \multicolumn{2}{c}{\pkg{partykit}, Monte Carlo} \\
 & & & & & \multicolumn{2}{c}{Bonferroni} \\
\cmidrule(lr){3-5} \cmidrule(l){6-7}
Observations & Predictors & Benchmark & No adjustment$^{\ast}$ & Max-type & 1 core & 32 cores \\ \midrule
500 & 20 & 9.6 & 0.16 & 0.44 & 12 & 0.7 \\
2,000 & 20 & 36 & 0.29 & 1.7 & 51 & 2.6 \\
2,000 & 100 & 39 & 0.57 & 1.8 & 74 & 4.0 \\
8,000 & 20 & 83 & 0.76 & 6.6 & 201 & 9.7 \\
8,000 & 100 & 97 & 1.7 & 6.2 & 308 & 15 \\
32,000 & 20 & 666 & 2.9 & 31 & 809 & 42 \\
32,000 & 100 & 665 & 6.2 & 19 & 1,221 & 77 \\
\bottomrule
\end{tabular}
\caption{Median fitting time in seconds over two repeats for 100-tree
classification forests on synthetic Gaussian predictors with weak signal in
the first five, on a 32-core machine. \citrees forests use
\code{n_jobs = -1}. Columns are as in Table~8 of the article.
$^{\ast}$A decomposition column, not a valid threshold test
(Section~\ref{sec:noadjust}).}
\label{tab:performance-scaling}
\end{table}

\section{The per-threshold test without adjustment} \label{sec:noadjust}

The no-adjustment rows and columns of Tables~7 and~8 test every candidate
threshold at the unadjusted level $\alpha_{\mathrm{B}}$. They show how much of
the fitting time the adjustment costs, but the resulting test does not control
its size. Table~\ref{tab:noadjust} reports its size in the Stage-B-only design
of Section~\ref{sec:threshold-design} at nominal 0.05 under the null. The size
is three to eight times the nominal level and grows with the number of
candidates. Adaptive and exhaustive stopping give the same rates. In the
complete-node design, where the Stage A test over five predictors precedes the
threshold test, the false-split rate without the adjustment stays between
0.0355 and 0.049 over the 36 design cells, so that design does not reveal the
failure.

\begin{table}[t!]
\centering
\small
\begin{tabular}{lrrrr}
\toprule
Task & Observations & 16 bins & 64 bins & 256 bins \\ \midrule
Classification & 100 & 0.156 & 0.218 & 0.226 \\
 & 200 & 0.152 & 0.242 & 0.298 \\
 & 500 & 0.172 & 0.270 & 0.314 \\
Regression & 100 & 0.210 & 0.282 & 0.286 \\
 & 200 & 0.230 & 0.298 & 0.324 \\
 & 500 & 0.236 & 0.334 & 0.384 \\
\bottomrule
\end{tabular}
\caption{Size at nominal 0.05 of the per-threshold test without the
Bonferroni adjustment over candidate thresholds, in the Stage-B-only design
under the null. Bins is the histogram setting.}
\label{tab:noadjust}
\end{table}

\section{Threshold test comparison} \label{sec:threshold-design}

The design and Table~\ref{tab:threshold-calibration} are those of the
companion article \citep{Milletich+Downes+Goley+Hirst:2026}, reported here
because they document the \code{threshold_test} argument. The comparison
measures the two Stage B tests under the multiple correlation and Pearson
correlation selectors, with adaptive and exhaustive stopping, in four parts.
Each part runs at five seeds, except the Stage-B-only design, which runs at
three. Both tests use the default \code{"auto"} budget, which exceeds the
floor at which adaptive stopping is inert. Candidate thresholds come from the
histogram method with 16, 64, and 256 bins; at 100 and 200 observations the
256-bin setting yields 100 and 200 candidates. Tables report medians over the
seeds, and a fit that exceeded its cap is reported as censored.

\subsection{Null calibration}

The first part estimates the complete-node false-split rate of a depth-one
tree under the global null, with five independent Gaussian predictors, an
independent response, sample sizes 100, 200, and 500, the three bin settings,
and 2,000 replicates per cell and seed. It measures the Stage A test over
five predictors followed by the Stage B test, so it lies below the Stage B
level alone. The Stage-B-only design isolates the threshold test with one
Gaussian predictor, the Stage A test disabled (\code{alpha_selector = 1}), a
depth-one tree, and nominal levels 0.02 to 0.50, at the same sample sizes and
bin settings. It uses 500 null and 200 alternative replicates per cell,
level, and seed, so that power can be compared at matched realized size.

The complete-node false-split rate is below 0.05 for both tests in all 36
design cells (Table~\ref{tab:threshold-calibration}). Under the per-threshold
Bonferroni test the rate falls as the bin setting grows, from 0.0196 at 16
bins to 0.0070 at 256 in classification with 500 observations and exhaustive
stopping. Adjacent thresholds induce nearly identical splits, and the union
bound charges for all of them. Under the max-type test the rate stays between
0.028 and 0.038 in every cell, with cell means of 0.032, 0.034, and 0.034 at
16, 64, and 256 bins. In the Stage-B-only design, the per-threshold size has
means 0.018, 0.011, and 0.006 at 16, 64, and 256 bins, and the max-type size
has means 0.033, 0.035, and 0.035. Every 95\% Clopper-Pearson upper limit is
below 0.056.

\begin{table}[t!]
\centering
\small
\begin{tabular}{lrrrrrr}
\toprule
 & & & \multicolumn{2}{c}{Per-threshold Bonferroni} & \multicolumn{2}{c}{Max-type} \\
\cmidrule(lr){4-5} \cmidrule(l){6-7}
Task & Observations & Bins & Adaptive & Exhaustive & Adaptive & Exhaustive \\ \midrule
Classification & 100 & 16 & 0.0182 & 0.0183 & 0.0342 & 0.0352 \\
 & 100 & 64 & 0.0127 & 0.0130 & 0.0353 & 0.0370 \\
 & 100 & 256 & 0.0091 & 0.0096 & 0.0360 & 0.0384 \\
 & 200 & 16 & 0.0200 & 0.0203 & 0.0326 & 0.0342 \\
 & 200 & 64 & 0.0144 & 0.0147 & 0.0333 & 0.0344 \\
 & 200 & 256 & 0.0086 & 0.0092 & 0.0329 & 0.0345 \\
 & 500 & 16 & 0.0196 & 0.0196 & 0.0295 & 0.0314 \\
 & 500 & 64 & 0.0142 & 0.0142 & 0.0322 & 0.0332 \\
 & 500 & 256 & 0.0072 & 0.0070 & 0.0319 & 0.0339 \\
Regression & 100 & 16 & 0.0272 & 0.0269 & 0.0315 & 0.0328 \\
 & 100 & 64 & 0.0174 & 0.0185 & 0.0337 & 0.0348 \\
 & 100 & 256 & 0.0138 & 0.0145 & 0.0335 & 0.0347 \\
 & 200 & 16 & 0.0282 & 0.0287 & 0.0323 & 0.0347 \\
 & 200 & 64 & 0.0234 & 0.0237 & 0.0351 & 0.0368 \\
 & 200 & 256 & 0.0124 & 0.0127 & 0.0356 & 0.0380 \\
 & 500 & 16 & 0.0262 & 0.0264 & 0.0286 & 0.0295 \\
 & 500 & 64 & 0.0198 & 0.0202 & 0.0288 & 0.0311 \\
 & 500 & 256 & 0.0106 & 0.0108 & 0.0285 & 0.0315 \\
\bottomrule
\end{tabular}
\caption{False-split rate of a depth-one tree under the global null at the
0.05 level, by Stage B test and stopping mode, over 10,000 replicates per
cell (five seeds of 2,000) with five independent Gaussian predictors and an
independent response. Bins is the histogram setting. Every 95\%
Clopper-Pearson interval lies below 0.05 (largest upper limit 0.042); the
widest half-width is below 0.004.}
\label{tab:threshold-calibration}
\end{table}

\subsection{Timing, real data, and power}

The second part times one tree on Gaussian predictors with a planted linear
signal at 1,000, 4,000, and 16,000 observations and the three bin settings,
as the median of three fits with a cap of 1,200 seconds per fit. With
exhaustive stopping, one classification tree at the 256-bin setting takes
3.5, 14, and 59 seconds under the per-threshold Bonferroni test, against 0.3,
1.4, and 6.3 seconds under the max-type test. The ratio grows with the bin
setting, as the cost orders of Section~2.2 predict. A per-threshold test held
at a fixed 999 permutations also costs of order $BK$, but it never splits at
64 or 256 bins, where $\alpha_{\mathrm{B}}/K$ falls below $1/1000$.

The third part fits one tree with the 256-bin setting and five-fold
cross-validation on six benchmark datasets fixed in advance by size: vowel,
spam, and page-blocks for classification and imports-85, residential, and
facebook for regression. The max-type test exceeded the per-threshold test's
held-out score (balanced accuracy or $R^2$) in 8 of 12 dataset and
stopping-mode cells and matched it in one. It was 0.9 to 8.3 times faster,
and its feature importances agreed with those of the per-threshold test at
Spearman correlations of 0.36 to 0.88. In the feature-ranking benchmark of
the companion article, which ranks 17 classification and 18 regression
feature-selection methods by mean rank across datasets, the max-type test
moved the forest from position 3.5 to 5 in classification and from 2 to 4 in
regression. Paired score differences favored the per-threshold test by at
most 0.0022 balanced accuracy and 0.009 $R^2$.

The fourth part estimates power with a depth-one tree on five Gaussian
predictors whose response depends on the first through a step at zero. The
step is a class-probability shift $\delta \in \{0.05, 0.10, 0.20\}$ in
classification and a mean shift $\delta \in \{0.2, 0.4, 0.8\}$ in regression,
at the same sample sizes and bin settings, with 500 replicates per cell and
seed. At a common nominal level, the max-type test splits more often, by 0.03
on average, and that gap comes from size. At matched realized size in the
Stage-B-only design, the difference in power, max-type minus per-threshold,
averages $-0.006$ over 108 cell and effect combinations (95\% bootstrap
interval over the 18 cells $-0.011$ to $-0.002$), with a bootstrap standard
error of about 0.017 per combination. One of 96 high-replicate combinations
lies beyond two standard errors, so per-cell differences are within Monte
Carlo error. The median distance of the chosen threshold from the planted cut
was 0.19 standard deviations under the max-type test and 0.15 under the
per-threshold test.

\section{NHANES inferential apparatus} \label{sec:nhanes-apparatus}

\subsection{Cohort and predictors}

Glycohemoglobin, fasting glucose, and the diabetes questionnaire define the
outcome and are excluded from the predictors. The 27 predictors are
demographic (age, sex, race and ethnicity, education, marital status, birth
in the United States, family income relative to the poverty threshold, and
the examination sampling weight), examination measurements (body mass index,
waist circumference, and systolic and diastolic blood pressure), laboratory
values (total and HDL cholesterol, triglycerides, creatinine, alanine
aminotransferase, uric acid, and gamma-glutamyl transferase), and
questionnaire items (ever smoked, alcohol frequency, sedentary minutes, sleep
hours, health insurance, and reported diagnoses of heart disease,
hypertension, and high cholesterol). The sampling weight is a survey-design
quantity retained as a natural high-cardinality predictor. The three permuted
columns are drawn once from the recorded seed in the primary run; their
cohort cardinalities are 3,464, 61, and 2. The \pkg{partykit} forest uses its
shipped defaults, including a candidate set of 6 predictors per node. Forests
are fitted with 100 trees on every available core.

\subsection{Test statistic and secondary contrasts}

Let $d_{ij}$ be the primary difference for fold $i$ of repeat $j$, let
$\bar{d}$ be its mean, let $s_d^2$ be the sample variance of the 50 fold
differences, and let $\bar{q}$ be the mean ratio of evaluation to training
sample size. Alongside the bootstrap, the lower-tailed corrected repeated
cross-validation statistic
%
\begin{equation}
  t =
  \frac{\bar{d}}
       {\sqrt{\left(1/50 + \bar{q}\right)s_d^2}}
  \label{eq:nhanes-corrected-cv}
\end{equation}
%
is referred to a $t$ distribution with 49 degrees of freedom
\citep{Bouckaert+Frank:2004}. The correction was derived for differences in
held-out performance, not for training-fold rank differences, so it is a
heuristic only. For the primary contrast, the ten repeat means ranged from
$-14.0$ to $-6.0$, and $t(49) = -4.62$, with a lower-tailed $p$-value below
0.001.

Two secondary contrasts use the same repeat-level bootstrap, with the
upper-tailed heuristic statistic adjusted across the pair by Holm's method.
The first is the Spearman correlation between a method's importances and the
columns' numbers of distinct training values, random forest minus \citrees.
It averaged 0.86 for the random forest, 0.36 for \citrees, and 0.23 for
\pkg{partykit}; the difference was 0.50, with repeat means from 0.44 to 0.53
and a 95\% bootstrap interval of 0.48 to 0.52. The second is the mean rank of
the eight predictors with at most three distinct values (sex, birth in the
United States, ever smoked, health insurance, marital status, and the three
reported diagnoses), random forest minus \citrees. It was 25.1 of 30 under
the random forest, 18.9 under \citrees, and 16.5 under \pkg{partykit}; the
difference was 6.19 ranks, with repeat means from 5.3 to 7.0 and a 95\%
bootstrap interval of 5.9 to 6.5. Both contrasts were in the specified
direction in all 50 folds and all ten repeats, and the Holm-adjusted
heuristic statistics, $t(49) = 15.8$ and $t(49) = 11.6$, both give $p$-values
below 0.001. Among real predictors the number of distinct values is
confounded with being a continuous laboratory measurement, so only the
permuted columns isolate cardinality from information. The two conditional
forests also ranked the binary items born in the United States and health
insurance in the lower half, so they separate informative from uninformative
binary items rather than promoting every low-cardinality column.

\subsection{Sanity checks and prediction}

Two checks were specified in advance as sanity conditions rather than
hypotheses. At least one of waist circumference, total cholesterol,
triglycerides, and HDL cholesterol should appear in each method's top 5 in a
majority of folds; all three primary methods did so in all 50 folds. Permuted
sex should fall outside each method's top 20 in a majority of folds; it did
in 50 of 50 folds for the random forest, 46 for \citrees, and 46 for
\pkg{partykit}.

For the descriptive prediction comparison, predictors are standardized on the
training fold, a ridge-penalized logistic regression is fitted to the 5 or 10
highest-ranked columns of each ranking, and the evaluation fold is scored by
log loss, area under the ROC curve, and Brier score. With 10 columns the log
loss was 0.319 for \pkg{partykit}, 0.320 for \citrees, and 0.333 for the
impurity-ranked random forest. With 5 columns the areas under the curve were
0.806, 0.792, and 0.776 for the same three rankings.

\subsection{Factor typing and redrawn controls}

The primary comparison passes the nominal columns to \pkg{partykit} as
numeric codes, although it natively tests and splits factors. A sensitivity
arm fitted in the same folds passes the ten nominal columns (sex, race and
ethnicity, marital status, birth in the United States, ever smoked, health
insurance, the three reported diagnoses, and permuted sex) as unordered
factors and keeps education and alcohol frequency as ordinal codes. Factor
typing leaves the cardinality pattern unchanged: the permuted sampling weight
has mean rank 26.6 against 26.4, the eight predictors with at most three
distinct values have mean rank 16.5 in both arms, and the Spearman
correlation between importance and number of distinct values is 0.20 against
0.23. The per-fold rank correlation with the numeric-coded ranking averages
0.91. The one large change is race and ethnicity, the only nominal column
with more than three categories, which moves from 21.6 to 12.9; marital
status moves from 23.6 to 21.7. The held-out area under the curve with 10
columns is 0.835 against 0.836.

In the out-of-bag control, the per-fold contrast of the permuted weight's
rank against \citrees is $+2.4$ ranks (standard deviation 6.1). In the
control without the Stage A test, the contrast is $+0.0$, and hypertension
moves from 5.8 to 11.1. A redraw run draws the three permuted columns
independently for each of the ten partition repeats. It moves the permuted
weight from 12.1 to 11.3 under the random forest and from 23.0 to 21.1 under
\citrees, and the primary contrast from $-11.0$ to $-9.8$ ranks.

Both forests rank by impurity decrease summed over splits, but not by the
same formula: \citrees sums node-level decreases, while \sklearn weights each
decrease by the node's share of the sample. The two forests also differ in
bootstrap and depth control, so the comparison does not isolate the split
search on its own.

\section{Number-to-producer map} \label{sec:producer-map}

Every reported number is produced by one component of the replication suite.
Table~\ref{tab:producer-map} lists the component and its output for each
table and study. Each component writes to a subdirectory of the output
directory named after the component, and each also writes a receipt with its
profile, seed, source revision, software versions, and input and output
hashes. Tracked copies of the published summaries are in the repository
directory \code{paper/results/tables}.

\begin{table}[t!]
\centering
\small
\begin{tabular}{@{}>{\raggedright\arraybackslash}p{3.7cm}>{\raggedright\arraybackslash}p{4.7cm}>{\raggedright\arraybackslash}p{7.2cm}@{}}
\toprule
Item & Component & Output \\ \midrule
Tutorial listings (Section 3) & \code{tutorial} & importances, paths, leaves, search summary, holdout predictions \\
Table 5 (calibration) & \code{calibration} & \code{selector_null_summary}, \code{root_null_summary} \\
Table 6 and cardinality rates & \code{calibration} & \code{cardinality_tree_trend_summary}, \code{cardinality_tree_method_summary}, \code{cardinality_forest_summary} \\
Matched comparison (Section 7.2, Section S2.2) & \code{behavior} & \code{behavior_summary}, \code{behavior_fold_summary} \\
Table 7, top panel & \code{head_to_head_timing}, \code{performance} & \code{head_to_head_summary.csv}, \code{performance_summary.csv} \\
Table 7, bottom panel & \code{performance_decomposition} & \code{performance_decomposition_summary.csv} \\
Table 8 and Table S1 & \code{head_to_head_timing} & \code{head_to_head_summary.csv} \\
Table 4 and Section S6 & \code{nhanes_diabetes} & \code{primary_inference}, \code{secondary_inference}, \code{rank_summary}, fold ranks and fold prediction metrics \\
Section S1 & \code{rdc_sensitivity} & \code{timing-summary.csv}, \code{performance-summary.csv}, \code{stability-summary.csv}, \code{comparison-summary.csv} \\
Section 7.4, Table S2, Section S5, Table S3 & \code{threshold_test} & null calibration, Stage-B-only design, real subset, and synthetic scaling summaries \\
\bottomrule
\end{tabular}
\caption{Replication component and output for each reported table and study.}
\label{tab:producer-map}
\end{table}

\bibliographystyle{jss}
\bibliography{refs}

\end{document}
